\newcommand{\figuraNguyen}[1]{
    \begin{figure}[H] % H fuerza la posición aquí
        \centering
        
        % --- FILA 1: EQL, EQL-DIV, DSR (Rutas conocidas) ---
        \begin{subfigure}[b]{0.5\textwidth}
            \centering
            % Ajusta la extensión (.png, .jpg, .pdf) si es necesario
            \includegraphics[width=\linewidth]{"Fotos Proyecto/EQL/Nguyen-#1"} 
            \caption{EQL}
        \end{subfigure}
        \hfill
        \begin{subfigure}[b]{0.5\textwidth}
            \centering
            \includegraphics[width=\linewidth]{"Fotos Proyecto/EQL-DIV/Nguyen-#1"}
            \caption{EQL-Div}
        \end{subfigure}
        \hfill
        \begin{subfigure}[b]{0.5\textwidth}
            \centering
            \includegraphics[width=\linewidth]{"Fotos Proyecto/DSR/Nguyen-#1"}
            \caption{DSR}
        \end{subfigure}
        
        \vspace{1em} % Espacio vertical entre filas
        
        % --- FILA 2: Phy-so, Ginn-Lp, Operon (Rutas asumidas) ---
        % NOTA: Debes verificar que estas carpetas existan o cambiar la ruta aquí
        %\begin{subfigure}[b]{0.5\textwidth}
        %    \centering
        %    \includegraphics[width=\linewidth]{"Fotos Proyecto/Phy-so/Nguyen-#1"}
        %    \caption{Phy-so}
        %\end{subfigure}
        %\hfill
        %\begin{subfigure}[b]{0.5\textwidth}
        %    \centering
        %    \includegraphics[width=\linewidth]{"Fotos Proyecto/Ginn-Lp/Nguyen-#1"}
        %    \caption{Ginn-Lp}
        %\end{subfigure}
        %\hfill
        %\begin{subfigure}[b]{0.5\textwidth}
        %    \centering
        %    \includegraphics[width=\linewidth]{"Fotos Proyecto/Operon/Nguyen-#1"}
        %    \caption{Operon}
        %\end{subfigure}
        
        \caption{Comparativa visual de modelos para el benchmark Nguyen-#1.}
        \label{fig:nguyen#1}
    \end{figure}
}


% --- INICIO DEL ANEXO ---
\appendix
\section{Resultados Visuales: Benchmarks Nguyen}

En esta sección se presentan las comparativas visuales para los casos seleccionados de Nguyen utilizando los seis modelos estudiados.

% Generación automática de las figuras
% Solo necesitas llamar al comando con el número correspondiente

\figuraNguyen{3}

\figuraNguyen{5}

\figuraNguyen{7}

\figuraNguyen{8}


\figuraNguyen{10}

\figuraNguyen{12}

\begin{figure}[H] % H fuerza la posición aquí
        \centering
        
        % --- FILA 1: EQL, EQL-DIV, DSR (Rutas conocidas) ---
        \begin{subfigure}[b]{0.5\textwidth}
            \centering
            % Ajusta la extensión (.png, .jpg, .pdf) si es necesario
            \includegraphics[width=\linewidth]{"Fotos Proyecto/EQL/Bat-25"} 
            \caption{EQL}
        \end{subfigure}
        \hfill
        \begin{subfigure}[b]{0.5\textwidth}
            \centering
            \includegraphics[width=\linewidth]{"Fotos Proyecto/EQL-DIV/Bat-25"}
            \caption{EQL-Div}
        \end{subfigure}
        \hfill
        \begin{subfigure}[b]{0.5\textwidth}
            \centering
            \includegraphics[width=\linewidth]{"Fotos Proyecto/DSR/Bat-25"}
            \caption{DSR}
        \end{subfigure}
        
        \vspace{1em} % Espacio vertical entre filas
        
        % --- FILA 2: Phy-so, Ginn-Lp, Operon (Rutas asumidas) ---
        % NOTA: Debes verificar que estas carpetas existan o cambiar la ruta aquí
        %\begin{subfigure}[b]{0.5\textwidth}
        %    \centering
        %    \includegraphics[width=\linewidth]{"Fotos Proyecto/Phy-so/Bat-25"}
        %    \caption{Phy-so}
        %\end{subfigure}
        %\hfill
        %\begin{subfigure}[b]{0.5\textwidth}
        %    \centering
        %    \includegraphics[width=\linewidth]{"Fotos Proyecto/Ginn-Lp/Bat-25"}
        %    \caption{Ginn-Lp}
        %\end{subfigure}
        %\hfill
        %\begin{subfigure}[b]{0.5\textwidth}
        %    \centering
        %    \includegraphics[width=\linewidth]{"Fotos Proyecto/Operon/Bat-25"}
        %    \caption{Operon}
        %\end{subfigure}
        
        \caption{Comparativa visual de modelos para las baterías de 25 celdas.}
        \label{fig:Bat25}
    \end{figure}
